\section{Ester và phản ứng ester hóa}

\subsection{Ethyl acetate và phản ứng ester hóa}

Khi đun ethylic alcohol với acetic acid có sulfuric acid đặc làm xúc tác, tạo ethyl acetate và nước:
\[
\ce{CH3COOH + C2H5OH <=>[H2SO4\ dac][t^\circ] CH3COOC2H5 + H2O}.
\]
Ethyl acetate là một ester, có mùi thơm đặc trưng và ít tan trong nước.

\begin{ghinho}{Điểm cần nhớ}
Phản ứng ester hóa là phản ứng thuận nghịch. Vì vậy không được mặc định toàn bộ acid và alcohol đều chuyển thành ester nếu đề không cho điều kiện hoặc hiệu suất phù hợp.
\end{ghinho}

\subsection{Thí nghiệm điều chế ethyl acetate}

Một quy trình điển hình gồm:
\[
\boxed{\ce{CH3COOH + C2H5OH + H2SO4}\ \text{đặc}}
\longrightarrow
\boxed{\text{đun nóng, hồi lưu}}
\]
\[
\boxed{\text{hỗn hợp sau phản ứng}}
\longrightarrow
\boxed{\text{thêm dung dịch \ce{NaCl} bão hòa}}
\longrightarrow
\boxed{\text{tách lớp}}.
\]

Trong thí nghiệm hồi lưu, ống sinh hàn giúp hơi ngưng tụ và quay trở lại bình phản ứng, hạn chế thất thoát các chất dễ bay hơi. Nước làm lạnh phải đi vào từ đầu dưới và ra ở đầu trên để áo nước luôn đầy.

Sau phản ứng, có thể thêm dung dịch \ce{NaCl} bão hòa để làm giảm độ tan của ester trong pha nước, giúp hai lớp phân tách rõ hơn. Có thể dùng phễu chiết để tách lớp ester.

\begin{notechem}
Khi đọc sơ đồ thí nghiệm, cần phân biệt rõ mục đích của từng thao tác: đun để tăng tốc độ phản ứng; sinh hàn để hạn chế mất chất; dung dịch \ce{NaCl} bão hòa để hỗ trợ tách lớp; phễu chiết để tách hai chất lỏng không trộn lẫn hoàn toàn.
\end{notechem}

\subsection{Tính toán trong phản ứng ester hóa}

\begin{vidu}
Cho \SI{12.0}{\gram} acetic acid phản ứng với \SI{11.5}{\gram} ethylic alcohol. Sau phản ứng thu được \SI{14.08}{\gram} ethyl acetate. Tính hiệu suất ester hóa. Cho $M_{\ce{CH3COOH}}=60$, $M_{\ce{C2H5OH}}=46$, $M_{\ce{CH3COOC2H5}}=88$ g/mol.
\end{vidu}
\begin{loigiai}
\[
n_{\ce{CH3COOH}}=\frac{12.0}{60}=\SI{0.200}{\mole},\qquad
n_{\ce{C2H5OH}}=\frac{11.5}{46}=\SI{0.250}{\mole}.
\]

\bangmol{7}
{\ce{CH3COOH} & + & \ce{C2H5OH} & \rightleftharpoons & \ce{CH3COOC2H5} & + & \ce{H2O}}
{0.200 & & 0.250 & & 0.200 & &}

Acetic acid là chất giới hạn. Khối lượng ester lí thuyết:
\[
m_{\text{ester,lt}}=0.200\times88=\SI{17.6}{\gram}.
\]
Hiệu suất:
\[
H=\frac{14.08}{17.6}\times100\%=80.0\%.
\]

\textbf{Vậy,} hiệu suất ester hóa là $80.0\%$.
\end{loigiai}

\begin{tuluyen}
Cho \SI{6.0}{\gram} acetic acid phản ứng với ethylic alcohol dư, hiệu suất ester hóa là $75\%$. Tính khối lượng ethyl acetate thu được.
\dapso{\SI{6.60}{\gram}.}
\begin{loigiaiBT}
\[
n_{\ce{CH3COOH}}=\frac{6.0}{60}=\SI{0.100}{\mole}.
\]

\bangmol{7}
{\ce{CH3COOH} & + & \ce{C2H5OH} & \rightleftharpoons & \ce{CH3COOC2H5} & + & \ce{H2O}}
{0.100 & & & & 0.100 & &}

\[
m_{\text{ester,lt}}=0.100\times88=\SI{8.80}{\gram},
\]
\[
m_{\text{ester}}=8.80\times0.75=\SI{6.60}{\gram}.
\]

\textbf{Vậy,} thu được \SI{6.60}{\gram} ethyl acetate.
\end{loigiaiBT}
\end{tuluyen}
